\documentclass[12pt,a4paper]{article}
\usepackage[utf8]{inputenc}
\usepackage[T1]{fontenc}
\usepackage[swedish]{babel} % Bra för svenska avstavningar
\renewcommand{\baselinestretch}{1.1}
\setlength{\parindent}{0pt}
\setlength{\hoffset}{-18pt}         
\setlength{\oddsidemargin}{0pt} % Marge gauche sur pages impaires
\setlength{\evensidemargin}{0pt} % Marge gauche sur pages paires
\setlength{\marginparwidth}{54pt} % Largeur de note dans la marge
\setlength{\textwidth}{481pt} % Largeur de la zone de texte (17cm)
\setlength{\voffset}{-18pt} % Bon pour DOS
\setlength{\marginparsep}{7pt} % Séparation de la marge
\setlength{\topmargin}{0pt} % Pas de marge en haut
\setlength{\headheight}{13pt} % Haut de page
\setlength{\headsep}{4pt} % Entre le haut de page et le texte
\setlength{\footskip}{27pt} % Bas de page + séparation
\setlength{\textheight}{720pt} % Hauteur de la zone de texte (25cm)

\begin{document}
\section{Boxall, P. \& J. Purcell (2022) Human resource management}
Human Resource Management (HRM) är en syftar att bygga, organisera och leda 
arbetskraften i en organisation för att uppnå organisatoriska mål. Kärnen i HRM ligger är följande:
\begin{itemize}
    \item \textbf{Ekonomiska mål}: Kostnadseffektivitet, flexibilitet, konkurrensfördelar
    \item \textbf{Sociopolitiska mål}: Social legitimitet och ledningsmakt
\end{itemize}
Boxall och Purcell (2022) Understryker att människor inte är resurser
i sig, utan oberoende aktörer som besitter mänskliga resurser såsom kunskap, färdigheter och energi. 
En effektiv HR-strategi kräver avvägningar och val som varierar beroende på:
\begin{itemize}
    \item \textbf{Industriell kontext}: Olika branscher har olika krav och utmaningar när det gäller HRM.
    \item \textbf{Nationella institutioner}: Lagstiftning, arbetsmarknadens struktur och kulturella normer påverkar HRM-praktiker.
    \item \textbf{Specifika personalgrupper}: Olika grupper av anställda kan ha olika behov och förväntningar, vilket kräver anpassade HR-strategier.
\end{itemize}

\subsection{Human resource management}
Processen där ledningen bygger upp en arbetskraft genom expansion och downsizing
\begin{itemize}
    \item \textbf{Expansion}: Rekrytering, urval, anställning och introduktion av nya medarbetare.
    \item \textbf{Downsizing}: Avveckling, avsked, uppsägning och avskedande av anställda.
\end{itemize}

\subsection{Mänskliga resurser kontra människor}
Organisationen är beroende av både mänskliga resurser såsom kunskaper och färdigheter.
\begin{itemize}
    \item \textbf{Mänskliga resurser}: De kunskaper, färdigheter och energier som människor besitter och som kan användas i arbetslivet.
    \item \textbf{Människor}: Oberoende aktörer som inte är resurser i sig, utan som har egna mål, behov och förväntningar.
\end{itemize}

\subsection{Human kapital och socialt kapital}
\begin{itemize}
    \item \textbf{Human kapital}: Mänsklig talang som är relevant för organisationens överlevnad
    \item \textbf{Socialt kapital}: Relationer och nätverk som skapar värde för organisationen
\end{itemize}

\subsection{De ekonomiska målen för HRM}
\begin{itemize}
    \item \textbf{Kostnadseffektivt arbete}: HRM måste bidra till ekonomisk lönsamhet genom att optimera arbetskraftens kostnader och produktivitet.
    \item \textbf{Organisatorisk flexibilitet}: HRM bör möjliggöra anpassning till förändringar i marknaden och arbetskraften, såsom teknologiska framsteg eller förändrade kundbehov.
\end{itemize}

\subsection{Human Resource Advantage (HR--fördelar)}
\begin{itemize}
    \item \textbf{Humankapitalfördel}: Anställa och behålla mer begåvade individer än konkurrenterna.
    \item \textbf{Organisatorisk processfördel}: Överlägsna sätt att kombinera individers talanger.
\end{itemize}

\subsection{Sociopolitiska mål för HRM}
Företag verkar i samhällen där lagar, normer och förväntningar finns
\\

Det finns tre pelare för institutioner:
\begin{enumerate}
    \item \textbf{Reglerande}: Lagstiftning kring anställning hälsa och säkerhet
    \item \textbf{Normativa}: Moralisk press att främja t.ex. mångfald och inkludering
    \item \textbf{Kulturellt-kognitiva}: Hur människor i en viss kultur tänker och agerar
\end{enumerate}

\subsection{Ledningsmakt (Management Power)}
    \begin{itemize}
        \item Ledningen behöver auktoritet för att samordna intressen och fatta beslut för organisationens bästa.
        \item \textbf{Maktsökande beteende}: Beslut som gynnar ledningens intressen snarare än organisationens övergripande mål.
    \end{itemize}

\subsection{Trade--offs}
\begin{itemize}
    \item \textbf{Arbetsmarknden}: Behov av att rekrytera talang kontra behovet av att kontrollera kostnader.
    \item \textbf{Kontroll vs samarbete}: Ledningen kan inte kontrollera allt utan är beroende av ansätlldas samarbete och vilja att bidra
    \item \textbf{Förändering}: Balasen mellan att ha en lojal personalstyrka och behovet av att minskar personalstyrkan vid nedgångar
    \item \textbf{Makt vs legitmitet}: Att maximera ledningens handlingsutrymme kontra behovet av att följa lagar och facklinga krav
\end{itemize}

\subsection{Särdrag för HR-strategier}
\begin{itemize}
    \item \textbf{Segmentering}: Olika moeller för olika grupper i.e "industriell modell med strikta regler", "professionell modell med autonomi" och "modellen för kärnkompetens med fokus på utveckling av talang"
    \item \textbf{Institutionell påverkan}: Natotionella institutioner påverkar strategin i.e Tyskland har starka system för yrkesutbildning som ger högkvalititativ arbetskraft, medan USA har mer flexibla arbetsmarknader som möjliggör snabb anpassning.
\end{itemize}






















\end{document}